\section{Visualisierungen}

Dieses Kapitel leitet die Anwendung her: zuerst die konkreten Aufgaben und die dafür nötigen
mentalen Modelle (\ref{sec:aufgaben}), daraus die Designanforderungen (\ref{sec:anforderungen}),
dann die drei Visualisierungen mit ihren Kodierungen und den verworfenen Alternativen
(\ref{sec:praesentation}) und schließlich die Interaktion, die sie zu einer Anwendung verbindet
(\ref{sec:interaktion}).

\subsection{Analyse der Anwendungsaufgaben}\label{sec:aufgaben}

Aus F1 bis F4 und den Bedürfnissen von Z1 bis Z3 ergeben sich sechs Aufgaben, bewusst als
Handlungen und nicht als Diagrammwünsche formuliert:

\begin{itemize}
  \item[\textbf{A1}] \emph{Orientieren.} Einen Einstiegszeitraum wählen, ohne vorher zu wissen, was
        ihn auszeichnet. (Voraussetzung aller weiteren Aufgaben)
  \item[\textbf{A2}] \emph{Rhythmen erkennen.} Feststellen, ob eine Schwankung dem Tagesgang, dem
        Wetterverlauf oder der Jahreszeit folgt. (F1, F3)
  \item[\textbf{A3}] \emph{Ausreißer lokalisieren.} Einzelne extreme Stunden in $8\,690$ Werten
        finden und ihren Zeitpunkt exakt ablesen. (F2)
  \item[\textbf{A4}] \emph{Ausdehnung prüfen.} Entscheiden, ob ein Extremwert eine Einzelstunde
        oder ein anhaltendes Muster ist. (F2)
  \item[\textbf{A5}] \emph{Tage vergleichen.} Mehrere Tage über alle relevanten Attribute
        gleichzeitig gegenüberstellen und ihre Ähnlichkeit beurteilen. (F3, F4)
  \item[\textbf{A6}] \emph{Gezielt auswählen.} Alle Tage finden, die zugleich mehrere
        Wertbedingungen erfüllen. (F2, F4)
\end{itemize}

\paragraph{Welche mentalen Modelle helfen dabei?} Die Aufgaben setzen voraus, dass die Nutzenden im
Kopf eine tragfähige Vorstellung des Systems aufbauen. Vier Modelle sind relevant, und jedes wird
durch eine der Ansichten gestützt.

\emph{M1: Die Zeit als zweidimensionales Raster.} Um A2 zu lösen, muss man Tagesgang und Jahresgang
als zwei unabhängige Achsen denken können und nicht als eine lange Linie. Erst wenn \enquote{jeder
Tag hat dieselben 24 Stunden} im Kopf steht, wird die Frage \enquote{tritt das immer zur selben
Tageszeit auf?} überhaupt stellbar. Für A2 und A4 ist das Modell notwendig: Ohne es bleibt eine
Häufung von Extremwerten eine unstrukturierte Menge von Zeitpunkten. Eine Darstellung, die die Zeit
in ein Tag\,$\times$\,Stunde-Raster faltet, vollzieht die Faltung bereits im Bild, statt sie dem
Betrachter aufzubürden.

\emph{M2: Die Deckungsbilanz.} Für A3 und A5 muss klar sein, dass der Erneuerbarenanteil ein
Verhältnis ist. Er kann steigen, weil die Erzeugung steigt oder weil die Last fällt. Ohne dieses
Modell werden Werte über $100\,\%$ als Fehler gelesen. Es entsteht, wenn Last, Solar, Wind und
Anteil gleichzeitig auf eigenen Achsen sichtbar sind.

\emph{M3: Der Preis als Folge der Knappheit.} A3 und A5 erfordern die Vorstellung, dass der Preis
kein weiteres Wettermerkmal ist, sondern aus der Deckungslücke folgt. Ob das zutrifft, ist selbst
eine empirische Frage. Die Anwendung muss den Zusammenhang deshalb sichtbar machen, statt ihn zu
unterstellen, und stellt Preis und Anteil gemeinsam dar, ohne sie zu verrechnen.

\emph{M4: Der Tag als Punkt in einem Merkmalsraum.} A5 und A6 sind ohne dieses Modell nicht lösbar:
\enquote{Ähnlichkeit zweier Tage} ist nur definiert, wenn ein Tag als Wertebündel gedacht wird und
Ähnlichkeit als Nähe in diesem Raum. Es ist das anspruchsvollste Modell und den Zielgruppen als
einziges unvertraut. Es braucht daher die stärkste visuelle Unterstützung: eine Darstellung, in der
ein Tag ein als Ganzes wahrnehmbares Objekt ist und die Achsen ihre ursprünglichen Einheiten
behalten.

Die Reihenfolge A1 $\rightarrow$ A2/A3 $\rightarrow$ A4/A5 $\rightarrow$ A6 entspricht dem Ablauf
\enquote{overview first, zoom and filter, details on demand} \cite{Shneiderman1996}.

\subsection{Anforderungen an die Visualisierungen}\label{sec:anforderungen}

\begin{itemize}
  \item[\textbf{R1}] \emph{Einstieg ohne Vorwissen.} Eine Ansicht muss das ganze Jahr auf einen
        Blick zeigen und als Filter dienen. (A1)
  \item[\textbf{R2}] \emph{Zwei Zeitskalen gleichzeitig.} Tagesgang und Jahresgang müssen in einer
        einzigen Darstellung sichtbar sein, ohne dass eine weggemittelt wird. (A2, M1)
  \item[\textbf{R3}] \emph{Vollständigkeit ohne Aggregation.} Alle $8\,690$ Stundenwerte müssen
        gleichzeitig darstellbar sein, damit seltene Ereignisse nicht durch Mittelung verschwinden
        ($441$ negative Preisstunden sind $5{,}1\,\%$ der Daten). (A3)
  \item[\textbf{R4}] \emph{Zusammenhängende Muster erkennbar.} Benachbarte Stunden müssen räumlich
        benachbart erscheinen, damit eine mehrstündige Phase eine Fläche und keine Punktwolke
        ergibt. (A4)
  \item[\textbf{R5}] \emph{Mehrere Attribute mit verschiedenen Einheiten.} Mindestens sieben
        Attribute (GW, \%, EUR/MWh, Anzahl) müssen für einen Tag gemeinsam lesbar sein, ohne auf
        eine gemeinsame Skala gezwungen zu werden. (A5, M2, M4)
  \item[\textbf{R6}] \emph{Wertbasierte Auswahl.} Wertebereiche auf mehreren Attributen müssen
        kombinierbar und die Treffermenge sichtbar sein. (A6)
  \item[\textbf{R7}] \emph{Kopplung.} Jede Beobachtung muss in mindestens einer anderen Ansicht
        überprüfbar sein; eine Auswahl darf nicht auf eine Ansicht beschränkt bleiben.
        (\cite{Roberts2007})
  \item[\textbf{R8}] \emph{Genaue Werte auf Abruf.} Da aus einer Farbe kein Zahlenwert ablesbar
        ist, muss jede Markierung ihre Werte auf Anfrage exakt ausgeben. (A3, A5)
  \item[\textbf{R9}] \emph{Ehrliche Kodierung.} Fehlende Daten dürfen nicht wie Nullwerte
        aussehen; Farbskalen müssen wahrnehmungsgerecht geordnet sein
        \cite{Borland2007,Harrower2003}; die wichtigste Größe gehört auf den genauesten Kanal
        \cite{Ware2012,Munzner2014}.
\end{itemize}

\begin{table}[H]
\centering\small
\caption{Zuordnung von Ansichten, Anforderungen und Aufgaben}\label{tab:anforderungen}
\begin{tabular}{lll}
\toprule
Ansicht & erfüllt vorrangig & unterstützt Aufgaben \\
\midrule
Zeitreihen-Übersicht (Vis 1) & R1, R7, R8 & A1, teilweise A2 \\
Stunden-Heatmap (Vis 2) & R2, R3, R4, R9 & A2, A3, A4 \\
Parallele Koordinaten (Vis 3) & R5, R6, R7 & A5, A6 \\
Detailpanel und Tooltips & R8 & A3, A5 \\
\bottomrule
\end{tabular}
\end{table}

\subsection{Präsentation der Visualisierungen}\label{sec:praesentation}

Die drei Techniken stammen aus drei verschiedenen der vorgegebenen Bereiche. Entscheidend ist aber
nicht diese formale Erfüllung, sondern dass jede diejenige Anforderung erfüllt, an der die beiden
anderen scheitern: Die Zeitreihe kann R3 nicht erfüllen, weil sie mittelt, die Heatmap nicht R5,
weil sie eine Größe zeigt, und die parallelen Koordinaten nicht R2, weil sie keine Zeitachse kennen.
Die Ansichten sind komplementär und nicht austauschbar. Bei jeder Technik werden Expressivität
(zeigt sie genau die Daten?) und Effektivität (nutzt sie für die wichtigste Größe den
wahrnehmungsstärksten Kanal?) im Sinne von \cite{Munzner2014,Ware2012} diskutiert.

\subsubsection{Visualisierung Eins: Zeitreihen-Übersicht der Monate}\label{sec:vis1}

\paragraph{Kodierung.} Die Ansicht zeigt zwölf Monatswerte. Der mittlere Erneuerbarenanteil wird
als Balkenhöhe kodiert (Kanal: Position auf gemeinsamer Achse bzw.\ Länge), der mittlere Preis als
Polylinie mit quadratischen Markern auf einer zweiten, rechts beschrifteten Achse in EUR/MWh.
Gitterlinien bei $0/25/50/75/100/125\,\%$ stützen das Ablesen. Farbe trägt keine Datenbedeutung,
sondern nur die Zuordnung zur Größe: Grün für den Anteil, Rot für den Preis, erklärt durch eine
Legende. Klick auf Balken oder Marker schaltet den Monat im Filter an bzw.\ aus, Hover liefert die
Monatswerte inklusive Solar- und Windanteil und Summe der negativen Preisstunden.

\paragraph{Warum diese Kodierung?} Der Erneuerbarenanteil ist die Leitgröße und liegt deshalb auf
dem genauesten Kanal (Position/Länge auf gemeinsamer Basislinie); der Preis als Nebengröße auf einer
Linie, die den Verlauf betont statt des exakten Wertes. Die doppelte Achse ist eine bewusst
eingegangene Schwäche: Der visuelle Abstand zwischen Balken und Linie hat keine Bedeutung, und die
scheinbare Korrelation hängt von der willkürlichen Skalierung der zweiten Achse ab. Sie ist
vertretbar, weil die Ansicht nicht zum Ablesen von Zusammenhängen dient, sondern als Navigations-
und Kontextebene für R1. Für Zusammenhänge sind Heatmap und parallele Koordinaten zuständig. Um aus
der Schwäche keine Fehlaussage werden zu lassen, ist die Preisachse explizit beschriftet und die
Interpretation im Bericht auf qualitative Aussagen beschränkt.

\paragraph{Alternativen.} \emph{Vollständige Stundenzeitreihe als Liniendiagramm.} Sie wäre
expressiver und erfüllte R3. Bei $8\,690$ Punkten auf einer Bildschirmbreite fallen jedoch rund
zwölf Werte auf ein Pixel; Überzeichnung macht Extremwerte ununterscheidbar von Rauschen, und der
Tagesgang erscheint als geschlossenes Band. Für A1 ist sie unbrauchbar. Ersatzlos gestrichen wird
sie deshalb nicht, sondern durch die Heatmap ersetzt, die dasselbe Datenvolumen ohne Überzeichnung
zeigt.

\emph{Horizon Graph.} Er löst genau dieses Auflösungsproblem: Durch Falten und Einfärben des
Wertebereichs erreicht er bei gleicher Höhe eine mehrfach höhere Datendichte und ist für lange
Zeitreihen Stand der Technik \cite{Aigner2011}. Verworfen, weil sein Lesen geübt sein muss (Bänder,
Vorzeichenfaltung) und er damit R1 verletzt, also den Einstieg ohne Vorwissen für Z2 und Z3. Für ein
reines Expertenpublikum wäre er die bessere Wahl.

\emph{Cycle Plot oder monatliche Boxplots.} Sie zeigten die innermonatliche Streuung, die der
Mittelwertbalken unterschlägt. Das ist ein realer Nachteil des gewählten Designs und wird in
Abschnitt~\ref{sec:anwendung1} sogar Gegenstand eines Anwendungsfalls. Verworfen, weil die Ansicht
zugleich Schaltfläche ist: Ein Balken ist eine große, eindeutige Klickfläche für den Monatsfilter,
eine Box mit Whiskern nicht. Die verlorene Streuungsinformation ist in der Heatmap vollständig
vorhanden, die Klickbarkeit wäre nicht ersetzbar.

\abb{zeitreihe}{0.92}{Zeitreihen-Übersicht: Balken zeigen den mittleren Erneuerbarenanteil je Monat,
die rote Linie den mittleren Preis auf der rechten Achse. Der Tooltip nennt die exakten Monatswerte.}
{Screenshot der Monatsübersicht mit Tooltip über einem Monatsbalken: zwölf grüne Balken, rote
Preislinie mit Markern, rechte Preisachse in EUR/MWh, Legende und Tooltip mit EE-Anteil, Preis,
Solar- und Windanteil sowie negativen Preisstunden.}

\subsubsection{Visualisierung Zwei: Pixelorientierte Stunden-Heatmap}\label{sec:vis2}

\paragraph{Kodierung.} Die Heatmap ist die dichteste Ansicht und setzt R2, R3 und R4 um. Jede Stunde
des Jahres ist eine Zelle: Die $x$-Achse läuft über die Tage des gefilterten Zeitraums, die
$y$-Achse über die $24$ Tagesstunden (0\,h oben, 23\,h unten), die Füllfarbe kodiert den
Erneuerbarenanteil. Ungefiltert sind das $366$ Spalten zu je $24$ Zellen; die Zellbreite wird
dynamisch als $\min(30,\,1500/\text{Tage})$ Pixel berechnet, sodass die Darstellung bei
Monatsfilterung von einer Pixel- in eine Kacheldarstellung übergeht, ohne dass sich die Semantik
ändert. Das ist eine Umsetzung des pixelorientierten Prinzips nach Keim \cite{Keim2000}: möglichst
viele Datenobjekte gleichzeitig, ein Objekt pro Pixel, mit einer Anordnung, die der Semantik der
Daten folgt.

\paragraph{Warum diese Anordnung?} Die Faltung der eindimensionalen Zeit in ein zweidimensionales
Raster ist der eigentliche Beitrag dieser Ansicht: Sie realisiert M1 direkt im Bild. Eine Struktur,
die vom Wetter kommt, erscheint als vertikaler Streifen über alle Stunden eines Tages; eine
Struktur, die vom Sonnenstand kommt, als horizontales Band in denselben Tagesstunden. Beide Fälle
sind ohne jede Rechnung unterscheidbar. Darin liegt die Antwort auf F1 und F3.

\paragraph{Farbe.} Die Skala ist sequentiell und interpoliert von einem dunklen Blauton (niedriger
Anteil) über Blaugrün zu hellem Grün (hoher Anteil) mit gemeinsam steigender Helligkeit. Sie folgt
damit der bekannten Konvention \enquote{Erneuerbare in Grün/Blau} und vermeidet bewusst eine
Regenbogenskala, deren nicht-monotone Helligkeit künstliche Grenzen erzeugt und die Ordnung der
Werte verschleiert \cite{Borland2007,Harrower2003}. Der Wertebereich wird bei $120\,\%$ begrenzt und
in der Legende als \enquote{$\geq 120\,\%$} ausgewiesen, sodass die $385$ Stunden über $100\,\%$
sichtbar bleiben, ohne den Kontrast des Hauptbereichs aufzubrauchen. Fehlende Stunden erhalten ein
neutrales, unbuntes Grau, das in der Farbskala nicht vorkommt und in der Legende eigens erklärt ist
(R9). Dass Farbe ein schwächerer Kanal als Position ist, wird bewusst in Kauf genommen und durch R8
aufgefangen: Der Tooltip liefert für jede Zelle Datum, Stunde, EE-Anteil, Preis, Last, Solar, Wind
und Nettohandel als Zahlen.

\paragraph{Alternativen.} \emph{Kalender-Heatmap nach van Wijk und van Selow} \cite{vanWijk1999}:
Sie ordnet Tage in einem Kalenderraster (Wochentage $\times$ Wochen) an und färbt jeden Tag nach
einer Kennzahl. Sie ist die naheliegendste Alternative und für Wochenrhythmen wie Werktag gegen
Wochenende überlegen. Verworfen, weil sie einen Tag auf einen Wert reduziert und damit R2 verletzt:
Der Tagesgang, die für Solar entscheidende Struktur, verschwindet. Ihre Stärke ist hier
zweitrangig, weil die Erzeugung im Gegensatz zur Last kaum Wochenstruktur hat.

\emph{Überlagerte Tageskurven (\enquote{Spaghettiplot}).} $366$ Kurven über $24$ Stunden zeigen den
Tagesgang exakt und erlauben das Ablesen von Werten. Verworfen wegen Überzeichnung: Im
Mittagsbereich ist die Dichte nicht mehr aufzulösen, und die für A3 zwingende Zuordnung einer Kurve
zu einem Datum ist ohne Interaktion unmöglich. Die Heatmap verzichtet auf exakte Werte (Farbe statt
Position) und gewinnt dafür die vollständige, überlagerungsfreie Zuordnung jeder Stunde zu ihrem
Ort im Bild.

\emph{Small Multiples je Monat.} Zwölf kleine Tag\,$\times$\,Stunde-Raster wären für den Vergleich
zwischen Monaten stärker, weil jeder Monat einen eigenen Rahmen bekommt. Verworfen, weil der
Übergang zwischen Monaten in getrennten Rahmen zerschnitten würde, etwa das kontinuierliche
Auseinanderlaufen des solaren Bandes von Februar bis Juni. Da die Anwendung den Monatsvergleich
bereits über Vis~1 und den Mehrfachfilter unterstützt, ist der durchgehende Zeitstrahl der größere
Gewinn.

\abb{heatmap}{0.98}{Stunden-Heatmap für das gesamte Jahr 2024: $x$ = Tag, $y$ = Stunde,
Farbe = Erneuerbarenanteil. Solares Mittagsband, windgetriebene vertikale Streifen und die grauen
Datenlücken im Mai und Oktober sind unmittelbar erkennbar.}
{Screenshot der Heatmap über alle zwölf Monate ohne Filter, mit Legende (0 bis 120\,\% plus Eintrag
für Datenlücken) und den Stundenbeschriftungen 0\,h, 6\,h, 12\,h, 18\,h, 23\,h an der linken Achse.}

\subsubsection{Visualisierung Drei: Parallele Koordinaten der Tagesprofile}\label{sec:vis3}

\paragraph{Kodierung.} Jeder gefilterte Tag ist eine Polylinie über sieben parallele, vertikale
Achsen: mittlere Last (GW), Solaranteil (\%), Windanteil (\%), Erneuerbarenanteil (\%), Nettohandel
(GW), mittlerer Preis (EUR/MWh) und Anzahl negativer Preisstunden. Jede Achse ist unabhängig auf
Minimum und Maximum im sichtbaren Datenbestand skaliert; beide Extremwerte stehen als Zahl an den
Achsenenden. Die Linien sind mit niedriger Deckkraft gezeichnet, sodass sich häufige Verläufe durch
Überlagerung von selbst verdichten. Hover hebt eine Linie rot hervor, Auswahl schwarz, vom Brushing
ausgeschlossene Tage werden fast vollständig ausgeblendet. Ist ein Achsenfilter aktiv, werden die
Treffer zusätzlich in kräftigem Grün gezeichnet: Ohne diese Hervorhebung bliebe die Treffermenge auf
derselben niedrigen Deckkraft wie das ungefilterte Bündel, und der Filter wäre zwar wirksam, aber
kaum sichtbar.

\paragraph{Warum diese Technik?} Parallele Koordinaten \cite{Inselberg1990,Heinrich2013} sind die
einzige der drei Techniken, die M4 aufbaut: Ein Tag wird zu einem einzigen Objekt, das man als
Ganzes verfolgen kann, und Ähnlichkeit zweier Tage wird als Ähnlichkeit zweier Linienverläufe
unmittelbar wahrnehmbar. Der entscheidende Vorteil gegenüber jeder Projektion ist, dass die Achsen
ihre Originaleinheiten behalten. Ein abgelesener Wert ist eine physikalische Größe und keine
Komponente. Genau das macht die Technik zur Grundlage für R6: Ein aufgezogener Bereich auf der
Preisachse ist eine Bedingung an den Preis und nicht an eine abstrakte Dimension.

\paragraph{Expressivität und Effektivität.} Expressiv sind parallele Koordinaten hier vollständig:
Alle sieben Attribute jedes Tages sind ohne Verlust dargestellt. Effektiv sind sie eingeschränkt.
Erstens hängt das Erkennen von Zusammenhängen von der Achsenreihenfolge ab. Nur benachbarte Achsen
zeigen ihre Korrelation als Bündelung bei gleichgerichtetem Verlauf oder als Kreuzung bei
gegenläufigem. Die Reihenfolge ist deshalb nicht beliebig, sondern folgt der Kausalkette des
Anwendungsgebiets: Last $\rightarrow$ Solar $\rightarrow$ Wind $\rightarrow$ EE $\rightarrow$
Handel $\rightarrow$ Preis $\rightarrow$ negative Stunden. Dadurch liegen die Paare nebeneinander,
deren Zusammenhang die Fragestellungen betreffen; das Paar (EE, Preis) zeigt das Auseinanderlaufen
der Linien, das der gemessenen Korrelation von $r = -0{,}80$ auf Tagesebene entspricht. Zweitens ist
bei $364$ Linien Überzeichnung unvermeidbar; sie wird durch niedrige Deckkraft in eine
Dichteinformation umgewandelt und durch Brushing beherrschbar gemacht.

\paragraph{Alternativen.} \emph{Streudiagramm-Matrix (SPLOM).} Sie zeigt jedes Attributpaar exakt
und ist für Korrelationen stärker. Verworfen, weil sieben Attribute $21$ Panels ergeben, jedes mit
$364$ Punkten: Ein einzelner Tag ist dann $21$-mal vorhanden und muss gedanklich wieder
zusammengesetzt werden. Genau diese Leistung soll M4 abnehmen.

\emph{Projektion (PCA oder MDS) mit Streudiagramm.} Sie löst das Überzeichnungsproblem elegant:
$364$ Tage als $364$ Punkte, Ähnlichkeit als Nähe, Cluster unmittelbar sichtbar. Verworfen, weil die
Achsen ihre Bedeutung verlieren. Auf die Frage \enquote{warum liegen diese Tage zusammen?} antwortet
eine Projektion nicht, und eine Bedingung im Sinne von R6 wie \enquote{zeige Tage mit negativem
Mittelpreis und hohem Solaranteil} ist in einem Komponentenraum nicht formulierbar.

\emph{Stern-Glyphen oder Radarplots je Tag.} Sie halten den Tag als Objekt zusammen und stützen M4.
Verworfen, weil der Vergleich von $364$ Glyphen ein serieller Suchvorgang ist, während parallele
Koordinaten alle Tage in einem gemeinsamen Bezugssystem übereinanderlegen. Hinzu kommt, dass Flächen
in Radarplots von der Achsenreihenfolge abhängen und systematisch fehlinterpretiert werden.

\abb{parallele-koordinaten}{0.98}{Parallele Koordinaten der Tagesprofile mit aktivem Brush auf der
Achse \enquote{Neg.\ Std.}: Nur Tage mit vielen negativen Preisstunden bleiben sichtbar, die
Statuszeile nennt die Trefferzahl.}
{Screenshot der parallelen Koordinaten über alle Tage, mit aufgezogenem Bereich im oberen Drittel der
rechten Achse, abgeblendeten übrigen Linien, Statuszeile \enquote{$n$ von 364 Tagen sichtbar} und
Schaltfläche zum Zurücksetzen der Achsenfilter.}

\subsection{Interaktion}\label{sec:interaktion}

Die drei Ansichten sind keine drei Diagramme, sondern eine Anwendung. Der gesamte
Interaktionszustand liegt in einem einzigen Elm-Model, aus dem alle Ansichten lesen; Konsistenz ist
dadurch nicht Konvention, sondern strukturell erzwungen. Vier Ebenen sind umgesetzt.

\paragraph{Ebene 1: Monatsfilter mit Mehrfachauswahl.} Monate lassen sich über eine
Schaltflächenleiste oder durch Klick auf einen Balken der Zeitreihe an- und abwählen; die Auswahl
ist eine Menge, sodass beliebige Kombinationen möglich sind (Juni bis August, oder Februar zusammen
mit November für einen Gegensatzvergleich). Sie bestimmt, welche Stunden in der Heatmap und welche
Tage in den parallelen Koordinaten und den Kennzahlenkarten erscheinen. \emph{Zweck:} A1 und
Reduktion der Datenmenge für A2 bis A6.

\paragraph{Ebene 2: Hover mit geteiltem Fokus und Tooltip.} Das Überfahren einer Heatmap-Zelle
oder einer Linie schreibt den Tag als \texttt{hoveredDay} in das Model. Beide Ansichten lesen diesen
Zustand, sodass die Hervorhebung immer gleichzeitig erscheint: Die Zellen des Tages erhalten einen
dunklen Rahmen, die Linie wird rot und in den Vordergrund gehoben, und das Detailpanel wechselt auf
diesen Tag. Zusätzlich erscheint am Mauszeiger ein Tooltip mit den exakten Werten. \emph{Zweck:} R7
und R8 ohne Zustandsänderung. Eine Vermutung lässt sich prüfen, ohne die bestehende Auswahl zu
verlieren.

\paragraph{Ebene 3: Mehrfachauswahl von Tagen.} Ein Klick auf Zelle oder Linie nimmt den Tag in
eine Liste auf, ein erneuter Klick entfernt ihn. Ausgewählte Tage bleiben in beiden Ansichten
dauerhaft schwarz hervorgehoben und erscheinen im Detailpanel als einzeln abwählbare Liste mit
EE-Anteil und Preis. Für den zuletzt berührten Tag zeigt das Panel sechs Kennzahlen und ein
Stundenprofil als kleines Balkendiagramm. Dessen Bezugslinien bei $50$ und $100\,\%$ und die
Stundenmarken alle sechs Stunden machen aus der reinen Verlaufsform eine ablesbare Größe. Die
$100$er-Linie ist dabei die inhaltlich wichtige, weil die Einspeisung oberhalb davon die Last
übersteigt.

Ab zwei ausgewählten Tagen tritt ein Vergleichsblock hinzu. Er besteht aus einer Tabelle mit allen
sieben Tagesattributen nebeneinander, einschließlich der Zahl negativer Preisstunden, die in der
Einzelansicht fehlt. Darunter stehen die Stundenprofile der ausgewählten Tage untereinander, sodass
sich die Tagesformen direkt gegenüberstehen (Solarbogen gegen flaches Windband). Beide sind nach
Datum sortiert und nicht nach Klickreihenfolge, damit derselbe Vergleich unabhängig von der
Bedienreihenfolge dasselbe Bild ergibt. \emph{Zweck:} A5, also der direkte Vergleich mehrerer Tage,
der mit einer Einfachauswahl nicht möglich ist.

\paragraph{Ebene 4: Achsen-Brushing.} Auf jeder der sieben Achsen lässt sich mit gedrückter
Maustaste ein Wertebereich aufziehen \cite{Becker1987}. Tage außerhalb werden abgeblendet und von
der Mausinteraktion ausgenommen. Brushes auf mehreren Achsen werden UND-verknüpft, sodass sich
zusammengesetzte Bedingungen formulieren lassen; eine Statuszeile nennt laufend die Trefferzahl, ein
Klick auf eine Achse löscht deren Filter, eine Schaltfläche alle.

Der Brush-Zustand wird dabei nicht in der Ansicht ausgewertet, sondern in \texttt{Main}. Von dort
speisen sich auch die Kennzahlenkarten, die sich damit auf die Treffermenge beziehen, und die
Heatmap, in der die Nichttreffer zurücktreten. Entscheidend ist, dass die Heatmap dafür abblendet
und nicht filtert. Alle Tage bleiben als Spalten stehen und die Zeitachse damit lückenlos.
Abgeblendete Tage sind von den grauen Zellen echter Datenlücken (R9) klar unterscheidbar. Ohne diese
Rückwirkung bliebe die stärkste Filteroperation der Anwendung auf eine einzige Ansicht beschränkt
und R7 wäre an dieser Stelle verletzt. \emph{Zweck:} R6, R7 und A6.

\begin{table}[H]
\centering\small
\caption{Kopplung der Ansichten: Jede Aktion wirkt über ihre Ursprungsansicht hinaus}\label{tab:kopplung}
\begin{tabular}{@{}p{3.9cm}p{2.9cm}p{3.6cm}p{3.6cm}@{}}
\toprule
Aktion & Zeitreihe & Heatmap & Parallele Koord. \\
\midrule
Monat wählen (Leiste oder Balken)
  & Balken markiert
  & zeigt nur diese Tage, Zellen werden breiter
  & zeigt nur diese Tage, Achsen reskalieren \\[0.3em]
Zelle überfahren (Heatmap)
  & unverändert
  & Tagesspalte gerahmt, Tooltip
  & Tageslinie rot hervorgehoben \\[0.3em]
Linie überfahren (Par.\ Koord.)
  & unverändert
  & Tagesspalte gerahmt
  & Linie rot, Tooltip mit 7 Werten \\[0.3em]
Zelle oder Linie klicken
  & unverändert
  & Tagesspalte dauerhaft markiert
  & Linie dauerhaft schwarz \\[0.3em]
Bereich auf Achse aufziehen
  & unverändert
  & Nichttreffer abgeblendet, Zeitachse bleibt vollständig
  & Treffer grün, Nichttreffer abgeblendet, Trefferzahl \\
\bottomrule
\end{tabular}
\end{table}

\paragraph{Was bewusst nicht umgesetzt wurde.} \emph{Umordnen der Achsen} wäre die wirksamste
Einzelverbesserung der parallelen Koordinaten \cite{Heinrich2013}, weil nur benachbarte Achsen auf
Korrelation prüfbar sind. Wir haben es zurückgestellt, weil die fachlich begründete feste
Reihenfolge die wichtigen Paare bereits benachbart macht und das Ziehen von Achsen mit dem Ziehen
von Brushes um dieselbe Geste konkurriert. \emph{Umschalten der Heatmap-Metrik} auf den Preis
bräuchte wegen dessen zweiseitiger Skala (negativ bis $936$\,EUR/MWh) eine divergierende Farbskala
mit Nullpunkt; zwei wechselnde Farbsemantiken in derselben Fläche sind eine bekannte Fehlerquelle
(siehe Abschnitt~\ref{sec:ausblick}). \emph{Ein Zeitraum-Brush in der Zeitreihe} wäre flexibler,
doch der Monat ist hier die natürliche Einheit, ein Klick billiger als ein gezogener Bereich, und
für tagesgenaues Eingrenzen steht die Heatmap zur Verfügung.

\abb{verknuepfung}{0.98}{Kopplung in Aktion: Monatsfilter Juli, der 7.\ Juli ist in der Heatmap
angeklickt und wird gleichzeitig in den parallelen Koordinaten hervorgehoben.}
{Screenshot mit Heatmap und parallelen Koordinaten untereinander, Monatsfilter Juli aktiv, die
Spalte des 7.\ Juli umrahmt, die zugehörige Linie hervorgehoben und ein Tooltip mit den
Stundenwerten.}
